%======================================================================
% Section 17: Conclusions and Scientific Outlook
%======================================================================

\section{Conclusions and Scientific Outlook}

\label{sec:conclusions}

\subsection{The shape of the field as this review has presented it}

\label{sec:17.1}

Read across the seventeen chapters of this review, generalized continuum mechanics for multiphysics wave propagation is a field with an unusually clear intellectual structure and an unusually weak evidential one. The intellectual structure is the work of a small number of decisive ideas, each of which the field has absorbed, extended, and preserved in a form that is recognisably the same. The evidential weakness is the condition under which those ideas have been tested.

The main modelling choices surveyed here are finite-rate heat-flux laws, alternative free-energy and entropy structures, additional field couplings, characteristic lengths, and time-history or fractional operators. Their mathematical consequences depend on the precise constitutive equations and parameter restrictions; a family name alone does not establish finite propagation, dissipation, stability, or well-posedness. Comparing these choices requires the same geometry, material data, boundary conditions, frequency convention, and validation target.

Important research directions include convergence results with error estimates, comparison of model predictions against measurements and independent implementations, and family-specific energy-flux derivations. Parameter identification from measurements that are independent of model fitting is also needed for meaningful comparison.

\subsection{Classical results and research questions}

\label{sec:17.2}

\subsubsection{Classical balance and model-specific calculations}

\begin{itemize}
\item The classical local total-energy density and outward-positive flux, from the displayed momentum and first-law balances, subject to symmetric Cauchy stress and no internal heat supply.
\item The high-wavenumber limits of the explicitly stated one-dimensional nonlocal dispersion relation in Section~\ref{sec:nonlocal}.
\item The stated-domain distinction between a quasistatic electric potential model and a full Maxwell system: the former has no displacement-current electromagnetic wave branch.
\end{itemize}

\subsubsection{Questions for future work}

\begin{itemize}
\item The explicit time-domain energy flux and dissipation function for each of the families listed in Section~\ref{sec:10.4}, obtainable in principle from the theory-specific constitutive evolution laws but not yet derived in a verifiable form.
\item The interface conditions for non-local and memory-dependent media, and the conditions under which the integral and differential formulations give consistent interface tractions.
\item Rigorous limit theorems between the theory families, with error estimates in a stated norm, which are the precondition for a quantitative rather than a qualitative cross-family comparison.
\item Experimental discrimination among the competing families on a common material, and identification of the relaxation, kernel, and length-scale parameters from data that were not generated by the theory being tested.
\item Reproducible numerical benchmarks executed independently, which would allow the validation axis of the taxonomy to be populated with something other than self-consistency.
\end{itemize}

\subsection{Methodological stance: tracing differences to assumptions}

\label{sec:17.3}

Differences between formulations may reflect constitutive choice, material symmetry, geometry, boundary conditions, harmonic convention, normalization, or approximation level; they should not be labelled contradictions without comparing these assumptions.

The stance has a cost, which should be stated. By declining to adjudicate between formulations, this review leaves the reader with a map rather than a recommendation in those places where the literature does not itself contain one. The map is, however, the more durable object: the particular theory that will be found correct in a given application is not yet decidable on the available evidence, and a review that pretended otherwise would be asserting more than the evidence supports. What the map does provide is a statement of what evidence would be needed to decide, and that is the more useful contribution at the present stage of the field.

See the single historical overview in Figure~\ref{fig:theory_evolution}.

\subsection{The three problems that matter most}

\label{sec:17.4}

Three questions are especially important for future analysis and experiment.

\paragraph{Energy accounting for the generalized families.}

Energy accounting for a generalized model should be derived from its constitutive equations and balance laws before modal ratios or absorption coefficients are compared. This is a model-specific derivation problem.

\paragraph{Limit behavior.}

Quantitative cross-family comparisons benefit from limit theorems with error estimates, rather than formal parameter substitutions alone.

\paragraph{Discriminating evidence.}

Discriminating experiments can compare pre-specified predictions for competing formulations against measurements on a common material. A common observable, measurement protocol, and model assumptions make such comparisons interpretable.

\subsection{Response to mock referee objections (WP6)}

\label{sec:17.6}

Six objections from the mock referee register (WP6-05, WP6-17, WP6-18, WP6-20, WP6-24, WP6-25) are answered here; each maps to ledger rows or to source-checked text.

\paragraph{WP6-05 --- Appendix table scale.} Appendix tables D, E, and F are source files that are not included in the PDF build; only appendices B and C are rendered (\texttt{main.tex}). No table overrun occurs in the 131-page PDF. Final-size render review is noted for any future reintroduction of appendix tables.

\paragraph{WP6-17 --- Equation audit scope.} The equation audit inventory (\texttt{check\_equation\_audit\_inventory.py}) covers every displayed equation group in Sections 3--14 and Appendix B: 137 active groups with one audit row each, plus 2 WITHHELD historical records with no active display. This extends the original Sections 4--10 scope.

\paragraph{WP6-18 --- Eq.~4.17 source anchor.} Eq.~\eqref{eq:4.17} (Roy Choudhuri 2007 TPL law) is source-anchored via OCL-169 (Class A) to Eqs.~(3), (4), and (9) of Roy Choudhuri (2007). The text retains only the first-order expansion with the stated lag-order condition $0\\leq\\tau_v<\\tau_T<\\tau_q$ and infers no common admissibility condition across variants.

\paragraph{WP6-20 --- Ledger line anchors.} Ledger line-anchor validation for rows such as OCL-097 and OCL-131 is scheduled for WP8; no rows were deleted from the ledger during this session.

\paragraph{WP6-24 --- GL and GN-I table cells.} Table T2 cells for Green--Lindsay and GN-I are marked ``Not established'' with Class B ledger rows OCL-481 and OCL-483. The table footnote states that ``Not established'' means the reviewed equations do not support a comparison in this audit, not that the property is absent. No family-wide GL/GN-I claim is asserted.

\paragraph{WP6-25 --- Piezoelectric constitutive displays.} The piezoelectric constitutive displays in Section~5 are source-anchored to Mindlin (1974), printed pp.~625--626, via OCL-195 (Class A). Sign and index conventions follow the cited source; no extension beyond the bounded classical linear thermopiezoelectric model is asserted.

\paragraph{W5.3 --- Captions state assumptions.} All figure captions state assumptions, scope, and illustrative status. Examples: sec01 (\emph{Parameter limits are model-specific}), sec05 (\emph{no universal mode count is asserted}), sec07 (\emph{correspond only to the particular one-dimensional laws stated}, \emph{Schematic, not a fit to data}), sec08 (\emph{Each step requires its own constitutive assumptions}), sec13 (\emph{The assumptions and validity domain must be stated}), sec16 (\emph{scope is limited to the displayed equations and assumptions}).

Against that assessment, the next decade of research in the field should pursue the following, in the order in which it would most improve the standing of the subject.

\begin{enumerate}
\item \textbf{Complete the energy flux derivations} for the generalized families, using the theory-specific constitutive evolution laws as the common template, and publish the dissipation function alongside the flux.
\item \textbf{Establish rigorous limit theorems} between the theory families, with explicit error estimates and in a stated norm.
\item \textbf{Formulate the interface conditions} for non-local and memory-dependent media, and demonstrate their consistency between the integral and differential representations.
\item \textbf{Create reproducible numerical benchmarks} that independent groups can execute using published data and code;
\item \textbf{Validate experimentally}, beginning with the thermal wave velocity and the interface energy partition, and with the dispersion measurement that distinguishes a non-local from a gradient mechanism.
\item \textbf{Extend the coupled-field models} to the materials for which the coupling is strongest---two-dimensional systems, metamaterials, and semiconductors---where the additional constitutive complexity is most likely to be repaid by predictive power.
\item \textbf{Characterize the parameters inversely}, so that the relaxation times, kernel parameters, and length scales are measured rather than fitted, and so that the number of free parameters in a model is matched to the number of independent measurements available.
\end{enumerate}
The next stage is to derive, prove, and measure the models already proposed, while identifying the regimes in which additional generalizations yield distinct, testable predictions. The review's organization around constitutive structures and shared analytical questions is intended to support such comparisons.

\subsection{Notation summary for Section 17}

\begin{table}[htbp]
\centering
\small
\caption{Notation introduced in Section 17.}
\label{tab:notation-sec17}
\begin{tabularx}{\textwidth}{@{}p{1.8cm}X p{2.2cm}p{2.1cm}@{}}
\toprule
\textbf{Symbol} & \textbf{Meaning} & \textbf{Units} & \textbf{Equation} \\
\midrule
$J_i$ & Total energy flux vector & W\,m$^{-2}$ & \eqref{eq:10.3} \\
$\sigma_{ij}$ & Cauchy stress tensor & Pa & Section~\ref{sec:foundations} \\
$q_i$ & Heat flux vector & W\,m$^{-2}$ & Section~\ref{sec:foundations} \\
$\tau_0$ & LS heat-flux relaxation time & s & Section~\ref{sec:thermoelastic} \\
$c_T$ & LS thermal characteristic speed & m\,s$^{-1}$ & Section~\ref{sec:thermoelastic} \\
$\omega$ & Angular frequency & rad\,s$^{-1}$ & Section~\ref{sec:methods} \\
$\Delta\bm{X}$ & Scattering matrix & --- & Section~\ref{sec:interfaces} \\
\bottomrule
\end{tabularx}
\end{table}
% LocalWords:
